\par
\SourceLeaf{391}%
Ferner sei \(r\). der Abstand des Mittelpunkts der gesuchten Ellipse vom Ursprung der Abszissenachse und \(\varphi\) der Winkel, den die vom Ursprung zum Mittelpunkt gezogene Gerade mit der Abszissenachse bildet. Die Koordinaten des Ellipsenmittelpunkts seien also \(r\cos\varphi,r\sin\varphi\). Daraus ergibt sich leicht für die Senkrechte von diesem Mittelpunkt auf die erste Seite des Vierecks die Länge
\[
=a-r\cos(A-\varphi)
\]
Die Senkrechten auf die übrigen drei Seiten lassen sich ebenso ausdrücken.

Bezeichnet man die große Halbachse der Ellipse mit \(\alpha\), die kleine Halbachse mit \(\beta\) und deren Neigung gegen die Abszissenachse mit \(\psi\), so ist offenbar \(A-\psi\) die Neigung der Senkrechten vom Mittelpunkt auf die erste Seite des Vierecks gegen die kleine Halbachse. Soll jene Seite die Ellipse berühren, so lässt sich diese Senkrechte aus bekannten Gründen durch
\[
\BareRoot{[\alpha\alpha\sin(A-\psi)^2+\beta\beta\cos(A-\psi)^2]}
\]
ausdrücken. Man hat also die Gleichung
\[
a-r\cos(A-\varphi)=\BareRoot{[\alpha\alpha\sin(A-\psi)^2+\beta\beta\cos(A-\psi)^2]}
\]
und ebenso drei weitere gleichartige Gleichungen, wenn man \(a\) und \(A\) durch die Zeichen für die übrigen Seiten ersetzt. Beseitigt man also die Wurzel und setzt der Kürze halber
\[
\begin{aligned}
rr-\alpha\alpha-\beta\beta&=t\\
\alpha\alpha-\beta\beta&=u
\end{aligned}
\]
so lauten unsere vier Gleichungen
\begingroup\fontsize{10}{12}\selectfont
\[
\begin{array}{rl}
\mathrm{I.}&2aa+t-4ar\cos(A-\varphi)+rr\cos2(A-\varphi)-u\cos2(A-\psi)=0\\
\mathrm{II.}&2a'a'+t-4a'r\cos(A'-\varphi)+rr\cos2(A'-\varphi)-u\cos2(A'-\psi)=0\\
\mathrm{III.}&2a''a''+t-4a''r\cos(A''-\varphi)+rr\cos2(A''-\varphi)-u\cos2(A''-\psi)=0\\
\mathrm{IV.}&2a'''a'''+t-4a'''r\cos(A'''-\varphi)+rr\cos2(A'''-\varphi)-u\cos2(A'''-\psi)=0
\end{array}
\]
\endgroup
Multipliziert man die erste Gleichung mit \(\sin2(A''-A')\), die zweite mit \(\sin2(A-A'')\), die dritte mit \(\sin2(A'-A)\) und addiert die Produkte, so erhält man (siehe Art.\,78 meiner \textit{Theoria motus corporum coelestium})
